\input{../def.tex}

\begin{document}

\avantgarde

\section{Life is awareness\,?}

Life is characterized by operating in theromdynamically open systems
where locally entropy is reduced,
while exporting the excess into the environment
so that the second law of thermodynamics is, of course, not violated.

This allows, for example,
a grasshopper to jump up from apparently completely resting,
while a rubber ball, for example,
in practice never starts to jump up by itself,
and if you drop it,
jumps up each time a little less high,
dissipating the directional energy into heat, random motion, increasing entropy.

{\color{xphi}What if it was precisely/identically the process of life described above
that would lead to a certain awareness,
as exemplified by creating directional motion or otherwise more order/structure locally\,?
%
In other words,
would life be the same as awareness,
as consciousness of varying degrees from cells to humans\,?}

This is tantalizing, at least an interesting hypothesis,
and see my posts in late January and early February 2026
for some circumstantial details,
especially via the book \textsl{Consciousness and the Universe}\raisebox{0.2ex}{\smaller*}$\!$,
edited by Roger Penrose et.\,al.

Note that this idea might also help regarding measurement in quantum mechanics,
as life and consciousness would come as one.
%
Would a Schrödinger’s cat already constantly measure itself by living\,?
%
Maybe path integrals would be interesting related to the idea,
since in the classical approximation they lead to a path with minimal action,
hence a directional, ordered path\,?

Despite not having a closed theory,
at least not yet,
one would assume that even with single cells with operating metabolisms
there would be no self-interference in two-slit experiments or similar setups,
simply because the cells would, like a Schrödinger’s cat, constantly measure themselves by living.

The main reason why I was susceptible to this new idea
were certainly my recent considerations about rest/move inside
since about mid 2024
that lead to images like a butterfly fluttering from flower to flower
or Odysseus from island to island,
and now to this idea \raisebox{0.5ex}{\texttildelow}\,28 Jan 2026 in the evening.

\vspace{9mm}
\noindent
{\footnotesize
\raisebox{0.1ex}{\smaller*} I looked inside the book on two subsequent evenings for the first time in my life
and (at least consciously) the most influence had the articles
\textsl{Does the Universe have Cosmological Memory? Does this imply Cosmic Consciousness?} (no year)
by Walter J.~Christensen Jr.,
mostly via the referenced article by Leo Szilard with original German title
\textsl{Über die Entropieverminderung in einem thermodynamischen System bei Eingriffen intelligenter Wesen} (1929),
and $\!$\textsl{The Stream of Consciousness} (1892) by William James,
while I skimmed most articles.
\par}

\end{document}
